% =====================================================================
% DocViz-Agent / QG-MDV — 논문 초안 v0.5 (2026-06-12)
% 변경점(v0.4.3 대비):
%  - framing 전환: "5축이 더 잘 생성한다" → "multi-document data construction이
%    본질 난제, rendering은 commoditized". headline은 source attribution gap.
%  - VSC/TMG는 headline contribution에서 제외 (VSC는 측정상 dead, TMG는 보조).
%  - method 기여를 cross-doc data construction + conflict adjudication 중심으로 재편
%    (Doc2Chart stage-1을 multi-doc으로 확장, proposition clustering 변형).
%  - 평가: 두 axis(visual: DiagramEval+VisJudge / data: Evidence F1 + attributed
%    value fidelity), attribution 비교의 fairness 설계(능력표+implicit proxy+
%    prompted-citation control) 명시.
%  - 3-source 실험 결과(node/path, Evidence F1) 표에 실측치 채움.
% 서술: 평범한 서술체 + English term 적극 사용.
% =====================================================================

\documentclass[11pt]{article}

\usepackage[hangul]{kotex}
\usepackage[english]{babel}
\usepackage{amsmath,amssymb}
\usepackage{booktabs}
\usepackage{multirow}
\usepackage{graphicx}
\usepackage{xcolor}
\usepackage{hyperref}
\usepackage{geometry}
\geometry{margin=1in}
\usepackage{enumitem}
\usepackage{tcolorbox}

\title{DocViz-Agent: Multi-Document 환경에서 \\
       Grounded Visualization을 위한 Data-Construction Agent}
\author{연구진 (placeholder)}
\date{2026년 6월 (v0.5)}

\begin{document}
\maketitle

\begin{abstract}
여러 문서에 흩어진 정보를 합쳐서 사용자 query에 맞는 visualization(chart 또는 diagram)을
만드는 일은 실제 업무에서 자주 필요하지만, 기존 연구는 대부분 single-document, single-chart에
머물러 있다. 이 논문은 이 과제를 Query-Grounded Multi-Document Visualization (QG-MDV)으로
정의한다. 우리가 던지는 핵심 질문은 "multi-document visualization에서 진짜 어려운 부분이
어디인가"이다. 실험을 해 보면, visualization을 그려내는 rendering 단계는 거의 모든 model이
비슷하게 해낸다. 즉 structure 품질을 재는 DiagramEval의 node/path alignment는 model 사이에
잘 갈리지 않는다. 반면 source attribution을 재면 격차가 크게 벌어진다. 그래서 우리는 이 과제의
본질적 난점이 rendering이 아니라 그 앞 단계, 즉 흩어지고 중복되고 충돌하고 단위가 제각각인
multi-document로부터 정확하고 traceable한 data를 구성하는 데 있다고 본다. 우리가 내세우는 중심 주장은 다음 한 문장이다. \emph{Multi-document visualization의 bottleneck은
rendering이 아니라, 여러 문서의 값을 reconcile하고 출처를 붙인 visual data specification을
구성하는 데 있다.} 그래서 우리는 최종 chart/diagram 자체보다, 그것을 만들기 위한 중간 산출물 —
\textbf{provenance를 보존한 Visual Data Specification (VDS)} — 을 논문의 핵심 산출물로 둔다.
DocViz-Agent는 cross-document에서 claim을 뽑아 visual element로 mapping하고, conflicting cell을
탐지/adjudicate하며, 각 visual element를 evidence에 binding하는 provenance-preserving pipeline
이다. 세 source(Loong, DocHop-QA, MultiHop-RAG)에서 제안 방법은 source attribution(Evidence F1)
에서 가장 강한 baseline을 큰 폭으로 앞선다(0.41 / 0.77 / 0.88 vs 약 0.00 / 0.14 / 0.04). 중요한
것은 이 grounding 개선이 visual quality를 희생하지 않고 얻어진다는 점이다(non-regression). 우리는
이 attribution 격차가 단순히 "우리만 citation을 출력해서" 생긴 측정 artifact가 아님을 보이기
위한 평가 설계(능력 분리, implicit proxy lower-bound, prompted-citation control)를 함께 제시한다.
\end{abstract}

% =====================================================================
\section{서론 (Introduction)}
% =====================================================================

LLM으로 chart나 diagram을 자동 생성하는 연구는 빠르게 발전해 왔다 \cite{han2023chartllama,
yang2024chartmimic}. 다만 이들은 대부분 하나의 표나 하나의 문서를 입력으로 받는다. 그런데
현실의 정보 작업은 여러 문서에 흩어진 내용을 통합한 뒤 시각화하는 경우가 많다. 예를 들어 금융
analyst는 여러 사업보고서에서 부문별 매출 추이를 모아 한 chart로 만들고, 의학 연구자는 여러
논문에서 약물 효과의 시간적 변화를 정리하며, policy analyst는 여러 정부 보고서 사이의 정책
변화를 diagram으로 그린다. 이런 작업은 입력이 multi-document이고, 사용자 의도가 자연어 query로
주어지며, 출력이 chart와 diagram을 모두 포함한다는 점에서 기존 single-document 연구가 다루지
못한 영역이다.

우리는 이 작업을 \textbf{Query-Grounded Multi-Document Visualization (QG-MDV)}으로 정의한다.
입력은 query $Q$와 document 묶음 $\mathcal{D}=\{D_1,\dots,D_n\}$ ($n\geq 2$)이고, 출력은 모든
visual element가 $\mathcal{D}$의 어느 문서에 근거하고 $Q$의 요구를 만족하는 visualization
산출물 $V$이다.

\textbf{이 논문의 출발 질문: 무엇이 진짜 어려운가.} 우리는 처음에 "잘 그리는 agent"를 만들면
이긴다고 가정했다. 그런데 실측을 해 보니 그렇지 않았다. visualization을 rendering한 뒤 structure를
재는 DiagramEval node/path alignment는 baseline과 제안 방법 사이에서 거의 갈리지 않았다(표
\ref{tab:nodepath}). rendering과 structure 생성은 요즘 model이면 대체로 해낸다는 뜻이다. 정작
크게 갈린 것은 source attribution이었다(표 \ref{tab:evidence}). 그래서 우리는 관점을 바꿨다.
\textbf{QG-MDV의 본질적 난점은 visualization을 그리는 단계가 아니라, 그리기 전에 multi-document
로부터 올바른 data를 구성하는 단계에 있다.} 구체적으로 다음을 풀어야 한다.

\begin{itemize}[leftmargin=*]
    \item \textbf{Cross-document reconciliation}: 같은 지표가 문서마다 다른 unit/scale/
    fiscal-year로 보고된다. 하나의 shared axis에 올리려면 통일해야 한다.
    \item \textbf{Conflict resolution}: 문서들이 같은 사실에 대해 서로 다른 값을 준다. 어떤
    값을 visual element로 encode할지 정하고 근거를 밝혀야 한다.
    \item \textbf{Completeness under distractors}: 핵심 정보가 여러 문서에 흩어져 있고 방해
    문서가 섞여 있다. 흩어진 series 구성원을 빠짐없이 모으고 무관한 것은 걸러야 한다.
    \item \textbf{Data-point attribution}: 생성된 visualization의 각 요소가 어느 문서, 어느
    근거에서 왔는지 추적 가능해야 한다.
\end{itemize}

이 네 가지는 multi-document QA/summarization 쪽(cross-document aggregation, 상충 정보, 중복,
정보 분산)과 document-to-chart 쪽(모호 수치의 값 추론, unit 변환, 큰 표의 부분 선택, attribution)
에서 각각 어려운 문제로 알려져 있다 \cite{tang2024multihoprag, wang2024loong, jain2025doc2chart,
zhao2024chartifytext}. QG-MDV는 이 둘이 겹치는 지점이다.

\textbf{기여 (Contributions).}
(C1) QG-MDV 과제 정의와, 학계 검증 multi-document benchmark에서 reasoning type 표지를 보존해
구축한 평가 corpus.
(C2) \textbf{provenance-preserving visual data construction pipeline}. cross-document claim을
visual element로 mapping하고, conflicting cell을 탐지하며, encode할 값을 adjudicate하고, 각
visual element를 evidence에 binding한다. 즉 "agent tool 집합"이 아니라, VDS(\S\ref{sec:vds})를
산출하는 구조화된 pipeline으로 contribution을 정의한다. 가장 가까운 선행 방법(Doc2Chart의
single-doc 추출/검증, proposition clustering의 요약용 fusion)을 그대로 쓰지 않고 multi-document
+ attribution 환경에 맞게 변형한다(\S\ref{sec:method}).
(C3) multi-document visualization에 맞춘 평가 체계. 기존 metric을 차용하되 lexical matching을
피하고, source attribution 비교의 fairness를 명시적으로 설계한다.
(C4) source attribution의 능력과 품질 격차가 backbone과 무관하게 일관됨을 보이는 진단적 발견.

% =====================================================================
\section{관련 연구 (Related Work)}
% =====================================================================

\textbf{Chart/diagram 생성.} ChartLlama \cite{han2023chartllama}, Text2Chart31
\cite{zadeh2024text2chart31}, ChartMimic \cite{yang2024chartmimic}은 single-table 또는
single-text에서 chart를 만든다. ChartMimic은 chart-to-code 평가에서 high-level(GPT-4V visual
similarity)과 low-level(text/layout/type/color) score를 쓴다. Doc2Chart \cite{jain2025doc2chart}
는 single-document에서 intent 기반으로 chart를 만들고, intent를 분해해 추출하고 정확성/완전성을
검증한 뒤 재추출하는 two-stage 방법을 쓴다. 우리는 Doc2Chart의 이 절차를 multi-document로
확장하고 conflict 처리와 provenance를 더한다(§4).

\textbf{Multi-document QA/summarization.} multi-document 작업의 어려움으로 cross-document
aggregation, 상충 정보 처리, 중복 제거, 정보 분산이 반복적으로 지적된다 \cite{ernst2022proposition,
tang2024multihoprag}. proposition-level clustering \cite{ernst2022proposition}은 문서들에서
명제를 뽑아 의미가 같은 것끼리 묶는다. 우리는 이걸 요약용 fusion이 아니라 visual series 조립용
으로 변형한다(§4).

\textbf{Visualization 평가.} DiagramEval \cite{liang2025diagrameval}은 diagram을 rendering한
뒤 VLM이 graph를 추출하고 node/path alignment로 평가하며, 사람 평가와 상관이 보고되어 있다.
VisJudge \cite{xie2025visjudge}는 visualization의 Fidelity-Expressiveness-Aesthetics를 재는
fine-tuned VLM judge로, 사람 일치도가 높다. Doc2Chart의 CHARTEVAL은 chart 값을 source 표로
역추적하는 attribution 기반 metric을 제안한다. 우리는 DiagramEval과 VisJudge를 visual axis로
차용하고, CHARTEVAL의 "값을 source로 추적한다"는 발상을 우리 source\_eid 기반으로 변형해
data axis에 쓴다(§5).

% =====================================================================
\section{과제 정의 (Task Formalization)}
% =====================================================================

QG-MDV의 입력은 query $Q$와 document 묶음 $\mathcal{D}=\{D_1,\dots,D_n\}$ ($n\geq 2$),
출력은 visualization 산출물 $V$이다. $V$는 chart 계열(Chart.js JSON)이나 diagram 계열
(Mermaid)로 표현되며, 모든 visual element가 $\mathcal{D}$의 특정 문서 chunk에 근거해야 한다.

각 query는 두 개의 독립 축으로 분류된다. reasoning 난이도(challenge type)는 multi-hop,
artifact planning, mixed artifact, distractor-heavy, contradiction의 다섯이고, output 구조
(output type)는 quantitative, relational, temporal, hierarchical, comparative의 다섯이다.
이 분류는 나중에 "어려운 type에서 우위가 큰가"를 보는 분석(§7)에 쓰인다.

\subsection{QG-MDV는 MD-QA나 document-to-chart와 무엇이 다른가}
\label{sec:diff}

reviewer가 가장 먼저 던질 질문은 "이거 multi-document QA에 chart rendering 붙인 것 아닌가"이다.
그렇지 않다. visualization 출력에는 평문 답에 없는 제약이 있고, 이 제약이 multi-document 환경과
만나면서 새 난점이 생긴다.

첫째, \textbf{shared axis 제약}이다. chart의 여러 막대/점은 하나의 축 위에 놓인다. 따라서 문서
마다 다른 unit/scale/fiscal-year로 보고된 값을 한 축에 올리려면 \emph{인코딩 전에} 공통 frame
으로 reconcile해야 한다. QA는 "문서 A는 \$4.5B, 문서 B는 4,500 million"이라고 둘 다 서술하면
되지만, 막대 하나는 단일 reconciled 값을 강제한다.

둘째, \textbf{visual element 단위 grounding}이다. QA의 citation은 문장(답)에 붙지만,
visualization에서는 막대 하나, 노드 하나, 엣지 하나가 각각 출처를 가져야 한다. 즉 attribution의
granularity가 문장이 아니라 visual element다.

셋째, \textbf{aggregated value의 provenance}이다. 막대 하나가 여러 문서의 값을 합/평균한
결과일 때, 그 하나의 visual element는 \emph{복수} 출처를 가리켜야 한다. 평문 답은 근거를 풀어
쓰면 되지만, visual element는 압축된 채로 다출처 provenance를 유지해야 한다.

넷째, \textbf{conflict가 single encoding을 강제}한다. 문서들이 한 값에서 충돌할 때, QA는 양측을
모두 서술할 수 있다. 그러나 막대 하나는 결국 한 값을 그려야 하므로, 충돌을 \emph{해소}하고 그
선택을 근거와 함께 남겨야 한다. 이 강제성이 conflict adjudication을 visualization-specific
문제로 만든다.

요약하면 QG-MDV는 MD-QA의 reasoning 위에 visualization 고유의 reconciliation·element-level
grounding·single-encoding 제약이 얹힌 과제다. 우리 method(\S\ref{sec:method})의 tool들은 바로
이 제약들을 직접 겨냥한다.

% =====================================================================
\section{방법 (Method): Multi-Document Data Construction Agent}
\label{sec:method}
% =====================================================================

본질 난점이 data construction이라면, 방법의 기여도 거기에 있어야 한다. DocViz-Agent는
다음 tool들로 구성된다. 핵심은 cross-document 추출/검증과 conflict adjudication 두 가지다.

\textbf{Cross-document iterative extraction \& validation.} Doc2Chart \cite{jain2025doc2chart}는
intent를 분해해 문서에서 관련 내용을 뽑고, 정확성/완전성/관련성을 검증한 뒤 재추출한다. 다만
이 검증은 single document 안에서만 이뤄진다. 우리는 이걸 \emph{여러 문서에 걸친} 추출로 넓히고,
검증 loop에 \emph{cross-document conflict 탐지}를 추가하며, 뽑은 값마다 source\_eid를 붙인다.
즉 Doc2Chart의 single-doc 검증을 multi-doc + provenance로 바꾼 것이 차별점이다.

\textbf{Claim clustering으로 visual series 조립.} proposition-level clustering
\cite{ernst2022proposition}은 명제를 뽑아 의미가 같은 것끼리 묶고 한 문장으로 fusion한다.
우리 목적은 요약이 아니라, 같은 (entity, metric, time)에 대한 claim들을 묶어 \emph{하나의
visual series 원소}로 mapping하는 것이다. cluster 안의 provenance는 유지하고 중복 series는
제거한다.

\textbf{Conflict adjudication.} multi-document 문헌은 상충을 난점으로만 지적하고, 이를
visualization 생성용으로 해소하는 방법은 없다. 우리는 같은 (entity, metric, time) cell에서
문서 간 값이 충돌하면 이를 탐지하고, policy(recency / source authority / majority)로 어떤 값을
encode할지 정한 뒤, 해당 visual element에 conflict flag와 provenance를 붙인다. 이건 신규
기여이고, contradiction type query로 직접 평가한다.

\textbf{Source-attributed output (SAO).} 각 data point / node / edge에 source\_eid를 붙인다.
이건 평가용 장치가 아니라 생성 단계에 들어 있는 mechanism이며, source attribution 결과의 토대다.

\subsection{Visual Data Specification (VDS)와 source\_eid binding}
\label{sec:vds}

위 tool들의 공통 산출물은 \textbf{Visual Data Specification (VDS)}, 즉 rendering 직전의
provenance 보존 중간 표현이다. VDS는 LLM이 자유 형식 DSL을 바로 쓰는 대신 거쳐 가는 구조화된
표현이며, 이것이 우리 contribution의 실제 산출물이다. chart의 경우 VDS는 각 datapoint를
$(\text{series}, \text{category}, \text{value}, \text{unit}, \text{time}, \text{source\_eid})$
tuple로 담고, diagram의 경우 node를 $(\text{id}, \text{label}, \text{source\_eid})$, edge를
$(\text{from}, \text{to}, \text{relation}, \text{source\_eid})$로 담는다. 여기서
$\text{source\_eid} = \{\text{doc\_id}\}\#\{\text{block\_id}\}\#\{\text{claim\_id}\}$이며, 한
visual element가 여러 출처를 합친 경우 source\_eid는 \emph{집합}이 된다(aggregated value의
다출처 provenance, \S\ref{sec:diff}).

\textbf{Binding이 rendering을 통과해 살아남게 하는 법.} VDS에서 Chart.js JSON / Mermaid로의
변환은 LLM 없이 deterministic rule로 한다. 변환할 때 각 VDS element의 source\_eid를 rendered
element의 index(예: \texttt{datasets[i].data[j]}, mermaid node id)에 정렬한 sidecar로 함께
출력한다. 따라서 attribution이 "문장에 붙은 인용"이 아니라 "rendering된 visual element에 binding된
provenance"로 보존된다. 이 deterministic 변환 덕에 값과 출처가 변환 과정에서 바뀌지 않는다.

\subsection{Conflict adjudication policy}
\label{sec:policy}

reviewer가 지적할 모호함 — recency/authority/majority가 충돌하면 무엇이 우선인가, 같은 단위·
시점에서만 비교하는가 — 을 정책으로 못박는다. 절차는 다음과 같다.

\begin{enumerate}[leftmargin=*]
    \item \textbf{Frame 정규화 후 비교}: 후보 값들을 공통 $(\text{entity}, \text{metric},
    \text{unit}, \text{time-granularity})$ frame으로 정규화한다. \emph{같은 frame에 속한 값들
    끼리만} 충돌로 본다. unit/time-granularity가 다르면 충돌이 아니라 별도 series이거나 정의된
    aggregation 대상이다(서로 다른 granularity를 majority로 투표하지 않는다).
    \item \textbf{우선순위}: 같은 frame 안에서 값이 갈리면 (1) source authority tier, (2)
    recency, (3) majority 순으로 적용한다. authority tier는 corpus 메타데이터(예: 1차/공식
    문서 $>$ 2차 인용)로 사전 정의하며, 정의 불가한 source는 동률로 둔다.
    \item \textbf{Tie와 unresolved}: 위로도 안 갈리면 한 값을 encode하되 \texttt{conflict}
    flag를 붙이고 충돌한 \emph{모든} source\_eid를 element에 함께 기록한다(은폐하지 않는다).
\end{enumerate}

이 정책 덕에 conflict adjudication은 prompt에 맡긴 판단이 아니라 재현 가능한 rule이 되며,
contradiction type query에서 직접 평가된다(\S\ref{sec:analysis}).

이전 버전(v0.4.3)에서 축으로 두었던 VSC(visual specification contract)는 측정상 효과가 거의
없어(deterministic to\_dsl이 이미 validity를 보장) headline에서 뺀다. 다만 그 deterministic
spec$\rightarrow$DSL 변환 자체는 \S\ref{sec:vds}의 source\_eid binding을 보존하는 장치로 재사용
한다. TMG(type-aware viz generation)도 보조로 내린다.

% =====================================================================
\section{평가 (Evaluation)}
% =====================================================================

\subsection{두 평가 axis}

평가는 두 axis로 구성하고, 둘 다 lexical(문자열) matching을 쓰지 않는다. lexical matching은
open-ended 생성에서 label 표현이 달라지면 깨지기 때문이다.

\textbf{Visual axis (rendering 기반).} 생성물을 실제로 rendering한 뒤 평가한다. DiagramEval
node/path alignment는 이미지에서 VLM이 graph를 뽑아 semantic(cosine) 정렬로 채점하므로 lexical이
아니다 — structure를 잰다. VisJudge는 rendering된 이미지를 Fidelity-Expressiveness-Aesthetics로
평가하는 외부 공개 judge다 — visual 품질을 잰다. 이 axis의 역할은 visual/structure 품질이 model
사이에서 거의 포화되어 있음을 보이는 것이다. 그래야 변별이 data axis에서 일어난다는 이야기가
성립한다. tree형 diagram에는 보조로 Tree Edit Distance를 쓴다.

\textbf{Data axis (source\_eid 기반).} 우리 기여가 data construction이므로 이 axis가 핵심이다.
Evidence F1은 attribution 정확도를 id 집합 연산으로 잰다. Attributed value fidelity는 각 data
point의 값이 \emph{그 point가 인용한 source}와 맞는지를 수치 tolerance로 본다. 후자는 CHARTEVAL의
"값을 source로 추적" 발상을 쓰되, 생성 표를 gold 표에 lexical로 정렬하는 원래 구현 대신
\emph{model이 붙인 source\_eid를 기준}으로 일치를 본다. 이러면 label lexical matching과 table
정렬 문제를 둘 다 피한다.

\subsection{Source attribution 비교의 fairness}

여기가 가장 민감한 부분이다. 제안 방법은 각 visual element에 source\_eid를 출력하지만,
baseline들(기존 visualization 생성 방법과 single-pass LLM)은 애초에 citation을 출력하는
mechanism이 없다. 그래서 id 집합 정확도를 그냥 비교하면 제안 방법이 자동으로 높게 나온다.
이걸 어떻게 다룰지가 문제다.

먼저, 이 비대칭이 곧 불공정은 아니다. multi-document에서 source를 attribute할 수 있다는 능력
자체가 우리 방법이 주는 정당한 기여다. 하지만 능력 차이만 점수로 내밀면 "너희만 id를 출력하게
만들었으니 당연하다"는 공격이 남는다. 그래서 우리는 한 점수로 합치지 않고 두 주장을 나눠서
보고한다.

\textbf{(1) 능력에 대한 구조적 주장.} 기존 방법과 single-pass LLM은 multi-document에서 visual
element를 source로 추적하는 mechanism을 산출 구조에 갖지 않는다. 이건 점수 비교가 아니라 방법의
구조적 속성이라, baseline에 인위로 citation을 시켜서 잴 게 아니다. "어떤 방법이 attribution을
산출 구조에 내장하는가"를 보이는 capability 표로 정직하게 제시한다(표 \ref{tab:capability}).
baseline에 citation을 강제하지 않는 이유는, 검색/정합 mechanism이 없는 single-pass model에
citation만 시키면 그럴듯하지만 틀린 인용을 만들어, 측정이 능력이 아니라 hallucinated citation을
재게 되기 때문이다.

\textbf{(2) 품질에 대한 정량 주장(방어 가능한 핵심).} 산출된 attribution이 실제로 얼마나
정확한가를 잰다. 제안 방법은 explicit attribution 정확도를 직접 잰다. baseline은 explicit
attribution이 없으니, visualization의 text 내용이 source chunk와 semantic하게 맞는지 보는
implicit matching을 \emph{관대한 lower-bound proxy}로 준다. 즉 "citation을 안 했어도 내용이
맞는 source에 닿아 있으면 인정"하는 가장 너그러운 조건을 baseline에 주는 것이고, 그래도 제안
방법이 앞서면 그 격차는 단지 id를 출력해서가 아니라 내용 자체가 맞는 source에 근거했기 때문이라고
말할 수 있다. 이 proxy가 공정하려면, 실제로 잘 근거한 baseline 산출물을 implicit matching이
제대로 인정하는지 사람 점검으로 확인해야 한다. 이 검증을 통과해야 lower-bound로 쓴다.

\textbf{(3) Prompted-citation control.} 위 두 주장을 보강하려고 통제 변형을 하나 둔다.
single-pass LLM에 data point마다 source를 citation하라고 지시한 변형이다. 목적은 baseline을
공정하게 만드는 게 아니라, citation을 출력하게 해도 cross-document 검색/정합 mechanism이 없으면
인용이 부정확함을 보이는 것이다. 이 변형이 explicit attribution 정확도에서 여전히 낮으면,
격차의 원인이 "id를 출력하느냐"가 아니라 "검색/정합/conflict 도구를 갖췄느냐"에 있음이 분리되어
입증된다. 이건 "attribution은 이미 흔한 기능 아니냐"는 공격에 대한 직접 방어이기도 하다.
multi-document visualization에서 attribution이 어려운 진짜 이유는, encode되는 값이 여러 문서에서
정합/집계되고 conflict가 해소된 결과이며, 문장이 아니라 visual element에 붙어야 하기 때문이다.

\subsection{차용 metric의 published score 활용}

DiagramEval 같은 차용 metric은 원 논문에 이미 score가 있는지 확인해서 두 가지로 쓴다. 하나는
metric 구현 검증이다. DiagramEval 원 논문이 자기 benchmark에서 보고한 node/path score를, 우리
구현을 그 공개 test set에 돌려 재현하면 채점 pipeline이 맞다는 걸 확인할 수 있다. 다른 하나는
external held-out에서의 baseline 재사용이다. 우리 자체 benchmark(QG-MDV)에는 어떤 선행 논문도
score를 보고한 적이 없고 baseline도 우리 과제로 이식한 적응판이라, 자체 benchmark에서는 baseline
측정을 생략할 수 없다. 다만 external held-out(Doc2Chart, Text2Vis 등)에서 그 benchmark의 원래
metric/protocol을 그대로 따르면, 원 논문의 baseline score를 인용하고 우리 방법만 새로 돌릴 수
있다. 이때도 우리가 재현한 한두 model의 score가 보고값과 비이상적으로 크게 어긋나지 않는지 먼저
확인하고, 크게 어긋나면 protocol 불일치를 의심해 인용하지 않는다.

% =====================================================================
\section{실험 (Experiments)}
% =====================================================================

\subsection{설정}

backbone은 Qwen3.5-397B(open frontier)를 주로 쓰고, closed backbone(GPT-5-mini, Claude Opus
4.8)은 API 예산이 풀리는 대로 추가한다. 비교 대상은 B1 MatPlotAgent, B2 nvAgent, B3 CoDA,
B4 ViviDoc(발표 SOTA 적응판), B5 Direct, B7 SelfRefine(단순 baseline), 그리고 B6
DocViz-Agent(제안)이다. metric은 DiagramEval node/path(seed-42), Evidence F1(explicit vs
implicit)이다. node/path는 temp 0.6에서 재실행 noise가 $\pm 0.04$ 수준이라 보조 신호로 보고,
최종본은 3-seed mean$\pm$std로 채울 예정이다.

\textbf{Fairness 통제.} 모든 arm은 동일한 document chunk 집합, 동일한 retrieval candidate,
동일한 context length를 받는다. 즉 B6가 더 좋은 정보를 더 많이 받아서 이기는 것이 아님을 보장
한다. 차이는 받은 정보가 아니라 그 정보로 무엇을 하는가(추출/정합/충돌해소/binding)에 있다.
한편 B6는 iterative extraction·validation·clustering·adjudication을 거치므로 single-pass보다
연산이 많다. 이 격차를 숨기지 않고, arm별 LLM call 수·평균 token·runtime을 \S\ref{sec:cost}에
보고한다(성능이 단지 더 많은 compute에서 온 것인지 판단할 수 있도록).

\subsection{Benchmark 구성과 gold의 신뢰성}
\label{sec:bench}

우리 corpus는 Loong, DocHop-QA, MultiHop-RAG 등 학계 검증 multi-document benchmark에서
출발한다. 이들은 원래 QA benchmark이므로, 우리는 (i) 원 benchmark의 document 묶음과 reasoning
type 표지를 보존하고, (ii) 각 query에 대해 gold visual data(plottable value 또는 graph
node/edge)와 그 값의 source attribution을 구성한다. gold attribution은 query를 만들 때 사용된
chunk의 provenance에서 결정론적으로 끌어오므로 채점기(VLM)와 분리된다. 다만 일부 gold reference
graph는 model로 생성되어 human-verified가 아니다. 이 한계 때문에 우리는 절대값보다 \emph{arm 간
상대 비교}에 의미를 두며(모든 arm이 동일 gold를 공유), gold attribution과 value의 정확성을
표본에 대해 사람이 audit하는 것을 평가의 일부로 둔다(\S\ref{sec:analysis}). 이 audit은 implicit
matching proxy의 공정성 검증(\S\ref{sec:fairness-result})과 함께 benchmark validity를 방어한다.

\subsection{Source attribution이 핵심 격차다 (Evidence F1)}

표 \ref{tab:evidence}가 핵심 결과다. source attribution(Evidence F1)에서 제안 방법은 세 source
전부에서 가장 강한 baseline을 큰 폭으로 앞선다. 이 격차는 node/path와 달리 noise 안에서 흔들리지
않고 일관된다.

\begin{table}[h]
\centering
\small
\begin{tabular}{lccc}
\toprule
arm & Loong & DocHop-QA & MultiHop-RAG \\
\midrule
\textbf{B6 (ours, explicit)} & \textbf{0.408} & \textbf{0.766} & \textbf{0.880} \\
B4 ViviDoc (implicit) & 0.004 & 0.141 & 0.041 \\
B7 SelfRefine (implicit) & --- & 0.131 & 0.012 \\
B5 Direct (implicit) & --- & 0.113 & 0.020 \\
B3 CoDA (implicit) & 0.000 & 0.058 & 0.008 \\
B2 nvAgent (implicit) & 0.000 & 0.041 & 0.000 \\
B1 MatPlotAgent (implicit) & 0.000 & 0.000 & 0.000 \\
\bottomrule
\end{tabular}
\caption{Source attribution (Evidence F1). B6는 explicit source\_eid, baseline은 implicit
embedding matching(cosine $\geq 0.75$). 세 source 전부에서 B6가 압도. 단 이 표의 해석은
\S\ref{sec:fairness-result}의 fairness 분석과 함께 읽어야 한다.}
\label{tab:evidence}
\end{table}

\subsection{Visual structure는 잘 갈리지 않는다 (node/path)}

표 \ref{tab:nodepath}는 DiagramEval node/path alignment다. path는 B6가 세 source 전부 최고지만,
node는 경쟁적이다 — Loong에서는 최고, DocHop에서는 B3와 동률, MultiHop에서는 B3에 뒤진다. 전반적
으로 model 사이 격차가 작다. 즉 rendering/structure 생성은 commoditized 되어 있고, 변별은 여기서
거의 일어나지 않는다. 이건 약점이 아니라 우리 주장의 근거다: 어려운 부분은 그리기가 아니라 data
construction이다.

\begin{table}[h]
\centering
\small
\begin{tabular}{l cc cc cc}
\toprule
& \multicolumn{2}{c}{Loong} & \multicolumn{2}{c}{DocHop-QA} & \multicolumn{2}{c}{MultiHop-RAG} \\
arm & node & path & node & path & node & path \\
\midrule
\textbf{B6 (ours)} & \textbf{0.396} & \textbf{0.340} & 0.432 & \textbf{0.552} & 0.241 & \textbf{0.191} \\
B5 Direct      & 0.388 & 0.309 & 0.401 & 0.354 & 0.227 & 0.087 \\
B2 nvAgent     & 0.359 & 0.320 & 0.393 & 0.281 & 0.237 & 0.085 \\
B7 SelfRefine  & 0.342 & 0.341 & 0.375 & 0.372 & 0.191 & 0.126 \\
B3 CoDA        & 0.321 & 0.247 & \textbf{0.437} & 0.343 & \textbf{0.289} & 0.129 \\
B4 ViviDoc     & 0.299 & 0.243 & 0.403 & 0.360 & 0.182 & 0.087 \\
B1 MatPlotAgent& 0.287 & 0.104 & 0.342 & 0.042 & 0.196 & 0.024 \\
\midrule
sample 수 & \multicolumn{2}{c}{100} & \multicolumn{2}{c}{100} & \multicolumn{2}{c}{96} \\
\bottomrule
\end{tabular}
\caption{DiagramEval node/path alignment F1 (seed-42). path는 B6가 세 source 전부 최고,
node는 경쟁적(MultiHop은 B3가 1위). model 간 격차가 작아 변별력이 약하다.}
\label{tab:nodepath}
\end{table}

발표 SOTA 적응판(B1--B4)이 multi-document에서 단순 Direct(B5)보다도 약한 경우가 많다는 점도
눈에 띈다. single-document chart에 특화된 방법이 multi-document로 오면 이점을 잃는다는 신호다.

\textbf{Non-regression claim.} 여기서 중요한 것은 우리가 "visual quality는 안 중요하다"고
말하는 게 아니라는 점이다. 우리 주장은 \emph{grounding을 크게 개선하면서 visual quality는
떨어뜨리지 않는다}는 것이다. node/path에서 B6가 최강 baseline과 동급 이상이라는 점이 그 1차
근거이고, 추가로 VisJudge(Fidelity-Expressiveness-Aesthetics)에서 B6가 baseline 대비 유의하게
나쁘지 않음을 보여 non-regression을 확정한다(표 \ref{tab:visjudge}, 측정 예정). 즉 data axis의
이득이 visual quality 희생 없이 얻어진다.

\begin{table}[h]
\centering
\small
\begin{tabular}{lccc}
\toprule
arm & Fidelity & Expressiveness & Aesthetics \\
\midrule
B6 (ours) & \_ & \_ & \_ \\
최강 baseline & \_ & \_ & \_ \\
\bottomrule
\end{tabular}
\caption{VisJudge 시각 품질 (측정 예정). non-regression — B6가 baseline 대비 유의하게 나쁘지
않음을 보인다. judge는 backbone pool과 분리한다.}
\label{tab:visjudge}
\end{table}

\subsection{Attribution capability와 fairness}
\label{sec:fairness-result}

표 \ref{tab:capability}는 \S5.2의 능력 주장을 정리한 것이다. 표 \ref{tab:evidence}의 큰 격차를
그대로 headline으로 쓰지 않고, (1) 능력 격차와 (2) attribution 품질을 나눠서 본다. baseline에는
implicit matching이라는 관대한 lower-bound를 줬는데도 거의 0에 가깝다는 점, 그리고 prompted-
citation control(아래)에서도 격차가 유지되는지가 핵심이다.

\begin{table}[h]
\centering
\small
\begin{tabular}{lcc}
\toprule
arm & explicit attribution 내장? & 측정 방식 \\
\midrule
B6 (ours) & 예 (SAO source\_eid) & explicit set F1 \\
B5+cite (control) & 지시로 출력 가능 & explicit set F1 \\
B1--B5, B7 & 아니오 (구조적 부재) & implicit embedding (lower-bound) \\
\bottomrule
\end{tabular}
\caption{Attribution capability 표. baseline은 산출 구조에 attribution mechanism이 없다.
이는 점수가 아니라 구조적 속성이며, 강제 citation으로 측정하지 않는다. 대신 prompted-citation
control(B5+cite)로 "id 출력만으로는 부족함"을 분리 측정한다.}
\label{tab:capability}
\end{table}

\textbf{Prompted-citation control (예정).} B5에 data point마다 citation을 지시한 B5+cite를
돌려 explicit Evidence F1을 잰다. 검색/정합 mechanism 없이 citation만 했을 때 정확도가 여전히
낮으면, 격차의 원인이 출력 형식이 아니라 agent의 data-construction 도구임이 분리 입증된다.

% =====================================================================
\section{분석 (Analysis)}
\label{sec:analysis}
% =====================================================================

\textbf{우위는 어려운 type에 집중되는가.} 우리 주장이 맞다면, 제안 방법의 우위는 data
construction이 어려운 challenge type(distractor-heavy, contradiction, multi-hop)에 몰려야 한다.
표 \ref{tab:bytype}는 세 source 결과를 challenge type별로 분해한 것이다(채울 예정). 쉬운 type
에서는 baseline도 비슷하고, 어려운 type에서 격차가 커지는 패턴을 기대한다.

\begin{table}[h]
\centering
\small
\begin{tabular}{lcc}
\toprule
challenge type & B6 Evidence F1 & 최강 baseline & $\Delta$ \\
\midrule
multi-hop & \_ & \_ & +\_ \\
artifact planning & \_ & \_ & +\_ \\
mixed artifact & \_ & \_ & +\_ \\
distractor-heavy & \_ & \_ & +\_ \\
contradiction & \_ & \_ & +\_ \\
\bottomrule
\end{tabular}
\caption{Challenge type별 분해 (채울 예정). 어려운 type에서 $\Delta$가 큰지로 "방법이 본질
난점을 푼다"를 검증한다.}
\label{tab:bytype}
\end{table}

\textbf{Tool ablation (예정).} cross-doc extraction 제거 / conflict adjudication 제거 / CIS
제거 / SAO 제거가 어려운 type 성능을 떨어뜨리는지 본다. ablation은 복구 산출이 아니라 native
산출 기반이어야 유효하다.

\textbf{Conflict adjudication (예정).} contradiction type query에서 arm별로 supported side를
맞게 encode하고 올바른 source에 attribute하는 정확도를 잰다. gold contradiction이 있는 sample에
대해 직접 측정한다. 이는 전체 Evidence F1보다 novelty가 더 살아나는 분석이므로 main paper에 둔다.

\subsection{Cost / latency / scalability}
\label{sec:cost}

multi-document agent는 single-pass보다 연산이 많으므로 "성능은 좋은데 너무 비싸지 않은가"라는
질문에 답해야 한다. 표 \ref{tab:cost}(측정 예정)에 arm별 평균 LLM call 수, 평균 token,
end-to-end runtime을 보고하고, document 수 증가에 따른 성능·비용 변화도 함께 본다. 목표는
B6의 이득이 단지 더 많은 compute에서 온 것이 아님을 보이는 것이다(예: B7 SelfRefine처럼 pass를
늘린 baseline과 비교).

\begin{table}[h]
\centering
\small
\begin{tabular}{lccc}
\toprule
arm & LLM calls & avg tokens & runtime (s) \\
\midrule
B5 Direct & \_ & \_ & \_ \\
B7 SelfRefine & \_ & \_ & \_ \\
B6 (ours) & \_ & \_ & \_ \\
\bottomrule
\end{tabular}
\caption{Compute budget (측정 예정). B6의 이득이 compute 증가만의 결과가 아님을 보인다.}
\label{tab:cost}
\end{table}

\subsection{Qualitative case study}
\label{sec:case}

task가 새롭기 때문에 숫자만으로는 독자가 감을 잡기 어렵다. 우리는 contradiction case 하나를
main figure로 제시한다(그림 \ref{fig:case}, 작성 예정). 구성은 다음과 같다: query, input
document들의 conflicting/distracting snippet, baseline visualization의 오류(예: 한 문서 값만
집어 잘못 encode하거나 출처 없음), DocViz-Agent의 중간 VDS(reconcile·adjudicate된 cell과
source\_eid), 최종 chart/diagram, 그리고 각 visual element의 attribution. 이 한 장이 §4의 tool
들이 실제로 무엇을 하는지, 왜 single-pass가 실패하는지를 직관적으로 보여준다.

\begin{figure}[h]
\centering
\fbox{\parbox{0.9\textwidth}{\centering \vspace{2em} [Contradiction case study: query →
conflicting snippets → baseline 오류 vs B6의 VDS(provenance) → 최종 viz → element별 attribution]
\vspace{2em}}}
\caption{Contradiction case study (작성 예정). 충돌하는 다문서 값에서 baseline은 임의 선택/무출처,
B6는 adjudication과 element-level attribution을 보인다.}
\label{fig:case}
\end{figure}

% =====================================================================
\section{논의와 한계 (Discussion and Limitations)}
% =====================================================================

\textbf{MD-QA와의 구분.} cross-document value reconciliation과 conflict adjudication은
visualization의 shared axis와 single-value 강제를 위한 연산이라, 평문 답을 내는 multi-document
QA에는 존재하지 않는다. prompted-citation control은 검색/정합 없이 citation만으로는 부족함을
보여, 우리 기여가 출력 형식이 아니라 data-construction mechanism임을 분리한다.

\textbf{한계.} (i) 출력 visualization 언어를 Chart.js와 Mermaid로 한정한다. (ii) visual 품질은
외부 judge(VisJudge)에 의존하므로 그 judge의 bias를 상속한다 — backbone pool과 분리해 순환을
막는다. (iii) multi-source attribution의 일부는 검색을 생략한 합성 근사로 산출됐으므로, 정식
주장 전에 검색 포함 산출로 재측정한다. (iv) gold reference graph가 model 생성이라 절대값보다
상대 비교에 의미를 두며 사람 검증으로 보강한다. (v) conflict adjudication 평가 sample이 제한적
이라 contradiction gold를 보강한다.

% =====================================================================
\section{결론 (Conclusion)}
% =====================================================================

이 논문은 QG-MDV 과제를 정의하고, 그 본질 난점이 rendering이 아니라 multi-document data
construction에 있음을 실측으로 보였다. DocViz-Agent는 cross-document 추출/검증과 conflict
adjudication을 중심으로 이 난점을 푸는 agent다. 세 source에서 source attribution 격차가 크고
일관되며, 우리는 이 격차가 측정 artifact가 아님을 보이는 fairness 설계(능력 분리, implicit
proxy, prompted-citation control)를 함께 제시한다.

% =====================================================================
\bibliographystyle{acl_natbib}
\begin{thebibliography}{99}

\bibitem{han2023chartllama}
Han et al. ChartLlama: A multimodal LLM for chart understanding and generation. arXiv:2311.16483, 2023.

\bibitem{yang2024chartmimic}
Yang et al. ChartMimic: Evaluating LMM's cross-modal reasoning capability via chart-to-code generation. ICLR 2025.

\bibitem{zadeh2024text2chart31}
Zadeh et al. Text2Chart31: Instruction tuning for chart generation with automatic feedback. EMNLP 2024 Main Oral.

\bibitem{jain2025doc2chart}
Jain et al. Doc2Chart: Intent-driven zero-shot chart generation from documents. EMNLP 2025 Main.

\bibitem{liang2025diagrameval}
Liang and You. DiagramEval: Evaluating LLM-generated diagrams via graphs. EMNLP 2025 Main.

\bibitem{xie2025visjudge}
Xie et al. VisJudge-Bench: Aesthetics and quality assessment of visualizations. arXiv:2510.22373, 2025.

\bibitem{ernst2022proposition}
Ernst et al. Proposition-level clustering for multi-document summarization. NAACL 2022.

\bibitem{wang2024loong}
Wang et al. Leave no document behind: Benchmarking long-context LLMs with extended multi-doc QA. EMNLP 2024 Oral.

\bibitem{tang2024multihoprag}
Tang and Yang. MultiHop-RAG: Benchmarking RAG for multi-hop queries. arXiv:2401.15391, 2024.

\bibitem{zhao2024chartifytext}
Zhao et al. ChartifyText: Automated chart generation from data-involved texts via LLM. arXiv:2410.14331, 2024.

\end{thebibliography}

\end{document}
